\section*{Choice assumptions and the original ultraproduct source}
\addcontentsline{toc}{section}{Choice assumptions and the original ultraproduct source}

\begin{proofboundary}{def:p03-ultrafilter-choice-boundary}{Boolean prime ideals, ultrafilters, and Choice}
Over \(\mathsf{ZF}\), the Boolean Prime Ideal Theorem is equivalent to the
Ultrafilter Lemma and is strictly weaker than the full Axiom of Choice, relative
to the usual consistency assumptions \cite{HowardRubin1998}.  Applied to the
cofinite filter on \(\mathbb N\), it suffices for the existence of a
nonprincipal ultrafilter.

The exact inherited development is written under the handbook's ambient
\(\mathsf{AC}\) convention.  Its assumption map is:
\begin{itemize}
  \item extending the cofinite filter to a nonprincipal ultrafilter needs only
        the Ultrafilter Lemma;
  \item the inherited type-space compactness proof invokes the already stated
        first-order compactness theorem in its ambient metatheory rather than
        claiming a new \(\mathsf{ZF}+\mathsf{BPI}\) calibration;
  \item the existential step in Los's theorem selects witnesses across the
        relevant index set, so the displayed proof retains full Choice;
  \item compactness via ultraproducts uses an ultrafilter on finite subtheories
        and also selects one model for each finite subtheory, so the displayed
        proof likewise retains full Choice; and
  \item the Keisler--Shelah theorem is used, under the ambient Choice
        convention, in the form Shelah published it; its hypotheses, and
        those of Keisler's earlier conditional theorems, are reconstructed in
        the next two paragraphs.
\end{itemize}

The original hypotheses of the cited ultrapower-isomorphism theorem are
these \cite[224--225]{Shelah1971Ultrapowers}.  An infinite cardinal
\(\lambda\) is fixed, and \(\mu\) is the least cardinal with
\(\lambda^\mu>\lambda\); the base of the exponent is \(\lambda\), not \(2\).
The index set of every ultrapower in the theorem is \(\lambda\) itself.  A
single ultrafilter \(D\) over \(\lambda\) is produced before any model is
considered, and it serves all models at once: for any two elementarily
equivalent models \(M\) and \(N\) of power less than \(\mu\), the ultrapowers
\(M^\lambda/D\) and \(N^\lambda/D\) are isomorphic.  The result is proved
without the generalized continuum hypothesis \cite[224]{Shelah1971Ultrapowers}.
The construction of \(D\) is a
transfinite recursion that begins by enumerating the pairs of models to be
handled and the \(\lambda\)-sequences from their universes in order type
\(2^\lambda\) \cite[230]{Shelah1971Ultrapowers}.  Such an enumeration is a
well-ordering of a set of functions on \(\lambda\), which the ambient Choice
supplies; no Choice-free version is claimed.  The full statement, the
cardinal computations, and the distinction between elementary equivalence,
isomorphic ultrapowers and isomorphism of the given models are in
\cref{sec:p03-w03-p03-018-shelah-ultrapowers}.

Keisler's earlier theorems have different hypotheses, read from his article
\cite{Keisler1961ElementaryClasses}.  The Axiom of Choice is assumed
throughout, with the well-ordering principle and the maximal principle
\cite[479]{Keisler1961ElementaryClasses}.  For the ultraproduct theorem, a
cardinal \(\alpha\) is fixed with \(\alpha\ge\omega\), with at least as many
elements as there are predicate symbols, and with
\(\varkappa(2^\alpha)=\alpha^+\), a single instance of the continuum
hypothesis (Keisler writes \(\varkappa(X)\) for the cardinality of \(X\) and
\(2^\alpha\) for the set of functions from \(\alpha\) to \(2\)).  The index set \(I\) has cardinality \(\alpha\), and every factor
has cardinality at most \(\varkappa(2^\alpha)\).  The conclusion supplies two
ultrafilters \(D\) and \(E\) on \(I\), each containing every subset of \(I\)
whose complement has fewer than \(\alpha\) elements, so neither is principal;
they are built after the factors are given and may differ
\cite[484--486]{Keisler1961ElementaryClasses}.  For ultrapowers the index set
is \(\alpha\) itself and the two structures have cardinality at most
\(\varkappa(2^\alpha)\) \cite[486]{Keisler1961ElementaryClasses}.  The
existential form, that elementarily equivalent structures have some
isomorphic ultrapowers, is stated under the full generalized continuum
hypothesis \cite[487]{Keisler1961ElementaryClasses}.  A single ultrafilter
\(D\) on \(\beta\), in a class \(Q(\beta^+)\) fixed before the structures,
appears only in results that the article announces without proof
\cite[490--492]{Keisler1961ElementaryClasses}; Shelah's summary of Keisler's
theorem \cite[224]{Shelah1971Ultrapowers} matches that announced form.  The
choices in Keisler's proof are the enumerations of the products of the
universes in order type \(\alpha^+\) by the well-ordering principle, a maximal
chain of pairs of sequences by the maximal principle, the choice of witnesses
at each step and in the disjoint-refinement lemma, and two applications of
the Ultrafilter Theorem, which is the Ultrafilter Lemma above.  These
hypotheses, the proofs from them, and the comparison with Shelah's are in
\cref{sec:p03-w03-p03-018-keisler}.


The distinctions are metatheoretic assumption boundaries, not new axioms of
the object theories.
\index{Boolean Prime Ideal Theorem}\index{Ultrafilter Lemma}\index{Axiom of Choice}
\end{proofboundary}

\begin{proofboundary}{def:p03-los-source-correction}{Original source of the ultraproduct theorem}
The ultraproduct construction and fundamental theorem occur in
\L o\'s's \emph{Quelques remarques, th\'eor\`emes et probl\`emes
sur les classes d\'efinissables d'alg\`ebres}, pages 98--113 of
\emph{Mathematical Interpretation of Formal Systems}
\cite{Los1955Ultraproducts}. The theorem applies to all first-order formulas;
the preceding assumption map specifies the Choice used by the proof given here.
\end{proofboundary}
